% CV.tex -- William's CV latex file.
\documentclass[10pt]{article}
%\usepackage{times}
%\usepackage{mathptm}
%\usepackage{revnum}
\newcommand{\comment}[1]{}
\newcommand{\myname}{William A.\ Stein}
\newcommand{\phone}{{\tt (206) 419-0925}}
\newcommand{\email}{{\tt wstein@sagemath.com}}
\newcommand{\www}{{\tt \url{https://wstein.org}}}
\newcommand{\address}{1212 East Barclay Court\\
Seattle, WA  98122}
\newcommand{\ptitle}{\em}
\usepackage{hyperref}
%\usepackage[hyperindex,pdfmark]{hyperref}

\hoffset=-.15\textwidth%
\textwidth=1.30\textwidth%
\voffset=-.05\textheight%
%\voffset=-.09\textheight%
\textheight=1.1\textheight%

\usepackage{fancyhdr,ifthen}
\pagestyle{fancy}
\cfoot{\thepage}  % no footers (in pagestyle fancy)
% running left heading
\lhead{\bfseries\LARGE\em \noindent{}\hspace{-.2em}\myname{}
        \hfill --\hfill Curriculum Vitae -- August 2026\vspace{-.2ex}\\}
% running right heading
\newcommand{\spc}{3.0em}
\rhead{\small
\phone
\hspace{\spc}$\cdot$\hspace{\spc}
\email
\hspace{\spc}$\cdot$\hspace{\spc}
\www}
% adjust, because the header is now taller than usual.
\setlength{\headheight}{7ex}

% bulleted list environment
\newenvironment{bulletlist}
   {
      \begin{list}
%         {$\bullet$}
         {$\cdot$}
         {
            \setlength{\itemsep}{.5ex}
            \setlength{\parsep}{0ex}
            \setlength{\leftmargin}{2.2em}
            \setlength{\parskip}{0ex}
            \setlength{\topsep}{0ex}
         }
   }
   {
      \end{list}
   }
%end newenvironment

\usepackage{revnum}

\newcounter{totalnumworkshops}
\setcounter{totalnumworkshops}{54}

\newcounter{totalnumworkshopsf}
\setcounter{totalnumworkshopsf}{8}

\newcounter{totalnumpapers}
\setcounter{totalnumpapers}{46}

\newcounter{totalnumgrants}
\setcounter{totalnumgrants}{21}

\newcounter{totalnumclasses}

\newcounter{totalnumstudents}
\setcounter{totalnumstudents}{7}

\newcounter{totalnumbooks}
\setcounter{totalnumbooks}{5}

\newenvironment{papers}
   {
   \begin{revnumerate}[\value{totalnumpapers}]
   {}
         {
            \setlength{\itemsep}{.5ex}
            \setlength{\parsep}{0ex}
            \setlength{\leftmargin}{0em}
            \setlength{\parskip}{0ex}
            \setlength{\topsep}{0ex}
         }
   }
   {
\end{revnumerate}
   }
%end newenvironment

\newenvironment{grants}
   {
   \begin{revnumerate}[\value{totalnumgrants}]
   {}
         {
            \setlength{\itemsep}{.5ex}
            \setlength{\parsep}{0ex}
            \setlength{\leftmargin}{0em}
            \setlength{\parskip}{0ex}
            \setlength{\topsep}{0ex}
         }
   }
   {
\end{revnumerate}
   }
%end newenvironment

\newenvironment{classes}
   {
   \begin{revnumerate}[\value{totalnumclasses}]
   {}
         {
            \setlength{\itemsep}{.5ex}
            \setlength{\parsep}{0ex}
            \setlength{\leftmargin}{0em}
            \setlength{\parskip}{0ex}
            \setlength{\topsep}{0ex}
         }
   }
   {
\end{revnumerate}
   }
%end newenvironment

\newenvironment{students}
   {
   \begin{revnumerate}[\value{totalnumstudents}]
   {}
         {
            \setlength{\itemsep}{.5ex}
            \setlength{\parsep}{0ex}
            \setlength{\leftmargin}{0em}
            \setlength{\parskip}{0ex}
            \setlength{\topsep}{0ex}
         }
   }
   {
\end{revnumerate}
   }
%end newenvironment

\newenvironment{books}
   {
   \begin{revnumerate}[\value{totalnumbooks}]
   {}
         {
            \setlength{\itemsep}{.5ex}
            \setlength{\parsep}{0ex}
            \setlength{\leftmargin}{0em}
            \setlength{\parskip}{0ex}
            \setlength{\topsep}{0ex}
         }
   }
   {
\end{revnumerate}
   }

\newenvironment{workshops}
   {
   \begin{revnumerate}[\value{totalnumworkshops}]
   {}
         {
            \setlength{\itemsep}{.5ex}
            \setlength{\parsep}{0ex}
            \setlength{\leftmargin}{0em}
            \setlength{\parskip}{0ex}
            \setlength{\topsep}{0ex}
         }
   }
   {
\end{revnumerate}
   }
%end newenvironment

\newenvironment{workshopsf}
   {
   \begin{revnumerate}[\value{totalnumworkshopsf}]
   {}
         {
            \setlength{\itemsep}{.5ex}
            \setlength{\parsep}{0ex}
            \setlength{\leftmargin}{0em}
            \setlength{\parskip}{0ex}
            \setlength{\topsep}{0ex}
         }
   }
   {
\end{revnumerate}
   }
%end newenvironment

\newenvironment{enumlist}
   {
   \begin{enumerate}
   {}
         {
            \setlength{\itemsep}{.5ex}
            \setlength{\parsep}{0ex}
            \setlength{\leftmargin}{0em}
            \setlength{\parskip}{0ex}
            \setlength{\topsep}{0ex}
         }
   }
   {
\end{enumerate}
   }
%end newenvironment

\newenvironment{dashlist}
   {
      \begin{list}
         {\hspace{2em}--}
         {
            \setlength{\itemsep}{.5ex}
            \setlength{\parsep}{0ex}
            \setlength{\leftmargin}{1em}
            \setlength{\parskip}{0ex}
            \setlength{\topsep}{0ex}
         }
   }
   {
      \end{list}
   }


% bold sans-serif label for margin labels.
%\newcommand{\marginlabel}[1]{{\large \textsf{\textbf{#1}}}}
\newcommand{\marginlabel}[1]{
\begin{minipage}[b]{0.8\labelwidth}{\large \textsf{\textbf{#1}}}\end{minipage}}

%\newcommand{\marginlabel}[1]{\textsf{\textbf{#1}}}
%\newcommand{\marginlabel}[1]{\textsf{\textbf{#1}}}
%\newcommand{\marginlabel}[1]{\textsc{\textsc{#1}}}

\newcommand{\entrylabel}[1]{\mbox{\marginlabel{#1}}\hfill}

\newcommand{\MainListlabel}[1]
   {
      \parbox[t]{\labelwidth}{\hspace{.8em}\marginlabel{#1}}
   }

% a list with fixed width marginlabels.
\newenvironment{MainList}[1]
   {
      \renewcommand{\entrylabel}{\MainListlabel}
      \begin{list}{}
      {
         \renewcommand{\makelabel}{\entrylabel}
         \setlength   {\itemindent}{-.65em}
         \setlength   {\labelwidth}{#1}
         \setlength   {\leftmargin}{\labelwidth}
         \setlength   {\itemsep}{3ex}
      }
   }
   {
      \end{list}
   }
%end newenvironment

\begin{document}
\begin{MainList}{88pt}
\item [Employment]
\begin{bulletlist}
\item {\bf SageMath, Inc.}: CEO/Founder, 2015--present.
\item {\bf University of Washington}: Prof. of
Math. (tenured), 2010--2019.
\item {\bf University of Washington}: Assoc. Prof. of
Math. (tenured), 2006--2010.
\item {\bf UC San Diego}: Assoc. Prof. of Math. (tenured),
2005--2006.
\item {\bf Harvard University}: Benjamin Peirce Asst. Prof. of Math.,
  2001--2005.
\item {\bf Harvard University}: NSF Postdoctoral Fellow, 2000--2004.
%\item Consultant for the Institute for Defense Analysis, 2002--present.
\end{bulletlist}

\item [Education]
\begin{bulletlist}
\item {\bf University of California at Berkeley},
      Ph.D. in Mathematics, 2000,
{\ptitle Explicit Approaches to Modular Abelian Varieties},
under H.~W.~Lenstra.
\item {\bf Northern Arizona University},
     B.S. in Mathematics, 1994.
\end{bulletlist}

\item [Prizes]
\begin{bulletlist}
\item {\bf Richard Dimick Jenks
Memorial Prize} for Excellence in
Software Engineering applied to Computer Algebra, 2013.
%\url{http://www.sigsam.org/awards/jenks/awardees/2013/}.
\item {\bf Trophe\'es du Libre}. The SageMath project won first
  prize in 2007 in Scientific Software (3000 euros, a laptop,
  books, server space, etc.).
\end{bulletlist}



\item [Grants]
\begin{grants}

%\item {\bf National Geospatial Agency}, grant to fund the workshop
%{\em Interactive Parallel Computation in Support of Research in Algebra, Geometry and Number Theory}, which I ran at MSRI Jan 29--Feb 2, 2007.
\item PI on {\bf NSA Grant}, {\em Sage: Open Source Math Software}, 2014--2017.
\item co-PI on {\bf NSF Grant}, DMS-7098841, {\em Sage-combinat: Developing and sharing open source software for algebraic combinatorics}, 2012--2015.
\item PI on {\bf NSF Grant}, DMS-7180474, {\em Explicit Approaches to Elliptic Curves and Modular Abelian Varieties}, 2012--2015.
\item co-PI on {\bf NSF Grant}, DMS-1062253, {\em REU: Inverse Problems for Electrical Networks}, 2011--2014.
\item co-PI on {\bf NSF Grant}, DMS-1020378, {\em Collaborative Research: UTMOST: Undergraduate Teaching in Mathematics with Open Software and Textbooks}, undergraduate curriculum development, 2011--2014.
\item PI on {\bf DOD Grant}, four Sage Bug-fixing workshops per year and other development support, 2011--2016.
\item PI on {\bf NSF Grant} DMS-1015114, {\em Sage: Unifying Mathematical Software for Scientists, Engineers, and Mathematicians}, four Sage Days workshops per year, 2011--2014.
\item UW Royalty Research Fellow (1-quarter teaching buyout), 2009--2010.
\item PI on {\bf NSF Grant} DMS-0821725, SCREMS grant for number
  theory, geometry, and software research (\$100K, purchased
  high end computers for UW).
\item PI on {\bf Google Grant}, for Sage development, Summer 2008 (amount: \$18K).
\item PI on {\bf Microsoft Grant}, port Sage to Windows, 2008--2009 (amount: \$32K).
\item PI on {\bf NSF Grant} DMS-0757627, {\em FRG: Collaborative Research: L-functions and Modular Forms}, 2008--2012
(amount: \$1.2 million).
\item co-PI on {\bf NSF Grant} DMS-0754486, {\em REU Site: Inverse Problems for Electrical Networks}, 2008--2011.
\item PI on {\bf NSF Grant} DMS-0713225, {\em SAGE: Software for Algebra and Geometry Experimentation}, 2007--2010.
\item PI on {\bf NSF Grant} DMS-0555776, {\em Explicit Approaches to the Birch and Swinnerton-Dyer Conjecture}, 2007--2010.
\item co-PI on {\bf NSF Grant} DMS-0602287, {\em Southwest Center for Arithmetic Geometry}, 2006--2009.
\item PI on {\bf NSF Grant,} DMS-0400386, {\em Explicit Approaches to Modular Forms and Modular Abelian Varieties}, 2004--2007.
\item PI on {\bf Sun Academic Education Grant} (\$70K Sun Fire V480 server), 2003.
\item From W.\thinspace{}R. Hearst III and Harvard
(\$20K for 12 Processor Cluster), 2002.
\item Clay Mathematics Institute Liftoff Fellowship, Summer 2000.
\item Berkeley Vice Chancellor Research Grant (6 Processor Cluster), 1999.
\end{grants}


\item[Publications] All papers are available at \url{https://wstein.org/papers/}.
 \begin{papers}

 \item {\em Torsion points on elliptic curves over number fields of small
 degree}, with Maarten Derickx, Sheldon Kamienny, and Michael Stoll, 2026
 preprint, \url{https://arxiv.org/abs/1707.00364}.

 \item {\em Beyond the black box}, with Jeroen Demeyer and Ursula Whitcher, 2016, Notices of the AMS.
 \item {\em Databases of elliptic curves ordered by height and distributions of Selmer groups and ranks} (22 pages), with Jennifer S. Balakrishnan, Wei Ho, Nathan Kaplan, Simon Spicer and Jamie Weigandt, 2016, in ANTS XII proceedings.

 \item {\em $p$-adic Heights of Heegner Points and Anticyclotomic Lambda-Adic Regulators} (34 pages), with Jennifer S. Balakrishnan and Mirela Ciperiani, 2013,
 Math. Comp.

\item {\ptitle A $p$-adic analogue of the conjecture of Birch and Swinnerton-Dyer for modular abelian varieties} (33 pages), with Jennifer S. Balakrishnan and J. Steffen M\"uller, 2013, Math. Comp.

\item {\ptitle A Database of Elliptic Curves over $\mathbf{Q}(\sqrt{5})$ -- First Report}
 (16 pages), with Jonathan Bober, Alyson Deines, Ariah Klages-Mundt, Benjamin LeVeque, R.~Andrew Ohana, Ashwath Rabindranath, Paul Sharaba, 2012, appeared in ANTS proceedings.

\item {\ptitle Numerical computation of Chow-Heegner points associated
to pairs of elliptic curves} (12 pages), 2012, to appear in Math. Comp. (as an appendix).

\item {\ptitle Sage: Creating a Viable Free Open Source Alternative
to Magma, Maple, Mathematica, and MATLAB} (9 pages), in the
FoCM 2011 proceedings.

\item {\ptitle Non-commutative Iwasawa theory for modular forms} (40 pages), with
John Coates, Tim Dokchitser, Zhibin Liang, Ramdorai Sujatha, 2013, in Proceedings of the London Math Society.

 \item {\ptitle Heegner Points and the Arithmetic of Elliptic Curves
     over Ring Class Extensions} (20 pages), with Robert Bradshaw,
   April 2012, J. Number Theory.

 \item {\ptitle Kolyvagin's Conjecture for Some Specific Higher Rank
     Elliptic Curves} (40 pages), 2011, submitted.

\item {\ptitle Computations About Tate-Shafarevich Groups Using
    Iwasawa Theory} (46 pages), with Christian Wuthrich, 2012, to
apear in Mathematics of Computation.

 \item {\ptitle The Sage Project: Unifying Free Mathematical Software
     to Create a Viable Alternative to Magma, Maple, Mathematica and
     MATLAB} (16 pages), 2010, in the Proceedings of the International
   Congress of Mathematical Software, Kobe, Japan.

\item {\ptitle Toward a Generalization of the Gross-Zagier Conjecture}
  (17 pages), 2010, Int.\ Math.\ Res.\ Notices.

\item {\ptitle Fast Computation of Hermite Normal Forms of Random
    Integer Matrices} (16 pages), with Clement Pernet,
Volume 130, Issue 7, July 2010, Pages 1675-1683, Journal of Number Theory.

\item {\ptitle Verification of the Birch and Swinnerton-Dyer
    Conjecture for Specific Elliptic Curves}, with G. Grigorov, A.
  Jorza, S. Patrikis, and C. Patrascu (26 pages), 2009, to appear in
  Mathematics of Computation.

\item {\em The Modular Degree, Congruence Primes and Multiplicity One}
  (16 pages), with Amod Agashe and Ken Ribet, 2012, in
a volume in honor of Serge Lang.

\item {\ptitle Explicit Heegner points: Kolyvagin's conjecture and
    non-trivial elements in the Shafarevich-Tate group}, with Dimitar
  Jetchev and Kristin Lauter (18 pages), 2008, Journal of Number
  Theory.

\item {\ptitle On the generation of the coefficient field of a newform
    by a single Hecke eigenvalue, with Koopa Koo and Gabor Wiese} (11
  pages), 2008, J. Th\'eor. Nombres Bordeaux.

\item {\ptitle Open Source Mathematical Software (opinion piece)} (1 page),
 with David Joyner, Notices of the AMS, November 2007.

\item {\ptitle Average Ranks of Elliptic Curve}, with
Baur Bektemirov, Barry Mazur and Mark Watkins (19 pages), May 2007,
Bulletins of the AMS.

\item {\ptitle Visibility of Mordell-Weil Groups} (20 pages), 2008,
  Documenta Mathematica.

\item {\ptitle Visualizing Elements of Shafarevich-Tate Groups at
    Higher Level}, with D. Jetchev (28 pages), 2008, Documenta
  Mathematica.

\item {\ptitle The Manin Constant}, with A. Agashe and K. Ribet (22
  pages), 2006, in the World Scientific Coates Memorial Volume.

\item {\ptitle Computation of $p$-Adic Heights and Log Convergence},
with B. Mazur and J. Tate (36 pages), 2006, in the Documenta
Mathematica Coates Memorial Volume.

\item {\ptitle SAGE: System for Algebra and Geometry Experimentation}
  with D. Joyner, (3 pages), in the SIGSAM Bulletin, 2005.

\item {\ptitle Modular Parametrizations of Neumann-Setzer
Elliptic Curves}, with M. Watkins, in IMRN 2004, no. 27, 1395--1405.

\item {\ptitle Studying the Birch and Swinnerton-Dyer Conjecture for
    Modular Abelian Varieties Using} MAGMA (23 pages), 2006, chapter in
  Springer-Verlag book edited by J.~Cannon and W.~Bosma.

\item {\ptitle Conjectures about Discriminants of Hecke Algebras of
Prime Level} (16 pages), with F.~Calegari, in ANTS VI, Vermont, 2004.
\item {\ptitle Constructing Elements in Shafarevich-Tate Groups of Modular
Motives},
   with N.~Dummigan and M.~Watkins, in ``Number theory and algebraic geometry---to Peter Swinnerton-Dyer on his 75th birthday'', Ed.
   M. Reid and A. Skorobogatov, pages 91--118.

 \item \emph{Approximation of eigenforms of infinite slope by
     eigenforms of finite slope}, with R.~Coleman, Geometric aspects
   of {D}work theory. {V}ol. {I}, {II}, Walter de Gruyter GmbH \&
   Co. KG, Berlin, 2004, pp.~437--449.

  \item {\ptitle $J_1(p)$ has connected fibers}, with B.~Conrad and
    B.~Edixhoven, Documenta Mathematica, {\bf 8} (2003), 331--408.
  \item {\ptitle Shafarevich-Tate Groups of Nonsquare Order},
  in  Progress in Math., {\bf 224} (2004), 277--289, Birkhauser.
  \item {\ptitle Visible Evidence for the Birch and Swinnerton-Dyer
Conjecture for Rank $0$ Modular Abelian Varieties} (30 pages),
with A.~Agashe, appeared in Mathematics of Computation.
  \item {\ptitle A Database of Elliptic Curves--First Report}
(10 pages) with M.~Watkins, in ANTS V proceedings, Sydney, Australia, 2002.
  \item {\ptitle Visibility of Shafarevich-Tate  Groups
    of Abelian Varieties}, with A.~Agashe,  J. Number Theory,  {\bf 97}
  (2002),  no. 1, 171--185.
  \item {\ptitle Cuspidal Modular Symbols are Transportable}, with H.~Verrill,
   LMS J.\ Comput.\ Math., {\bf 4} (2001), 170--181.
\item Appendix to Lario and Schoof's {\ptitle Some computations
    with Hecke rings and deformation rings}, with A.~Agashe,
     Experiment. Math. {\bf 11} (2002), no.~2, 303--311.
  \item {\ptitle There are genus one curves over $\mathbf{Q}$
         of every odd index}, J. Reine Angew. Math. {\bf 547}
(2002), 139--147.
  \item {\ptitle Component groups of purely toric quotients of
semistable Jacobians}, with B.~Conrad, Math. Res. Lett.,  {\bf
8}  (2001),  no. 5--6, 745--766.
  \item {\ptitle The field generated by the points of small prime
       order on an elliptic curve}, with L.~Merel,
       Int.\ Math.\ Res.\ Notices,  2001, no.~20, 1075--1082.
  \item {\ptitle An introduction to computing modular forms using
modular symbols} (12 pages), in MSRI Publications (Volume 44),
Algorithmic Number Theory: Lattices, Number Fields, Curves
and Cryptography, Cambridge University Press, 2008.
  \item {\ptitle A mod five approach to modularity of icosahedral Galois
         representations}, with K.~Buzzard, Pac. J.
Math.,  {\bf 203}  (2002),  no. 2, 265--282.%
  \item {\ptitle Lectures on Serre's conjectures},
     with K.\thinspace{}A.~Ribet,
     in Arithmetic Algebraic Geometry, IAS/Park City Math.
     Inst. Series, Vol.~9, 143--232.
  \item {\ptitle Component groups of quotients of $J_0(N)$},
     with D.~Kohel,   Proceedings of the 4th International
     Symposium (ANTS-IV), 2000, 405--412.

\item {\ptitle Empirical evidence for
     the Birch and Swinnerton-Dyer conjectures for modular Jacobians of
     genus 2 curves},
     with E.\thinspace{}V.~Flynn, F.~Lepr\'{e}vost,
     E.\thinspace{}F. Schaefer,
     M.~Stoll, J.\thinspace{}L.~Wetherell,
     Math.\  of Comp. {\bf 70} (2001), no. 236, 1675--1697.


%  \item {\ptitle Fallacies, Flaws, and
%         Flimflam \#92:   An Inductive Fallacy} (1 page),
%     with A.~Riskin,  College Math.\ J.\ 26:5 (1995), 382.
%\newcounter{local}\setcounter{local}{\value{enumi}}
 \end{papers}



\item [Books]
\begin{books}
%\setcounter{enumi}{\value{local}}
\item {\ptitle Prime Numbers and the Riemann Hypothesis} (154 pages), with
  B. Mazur (see \url{https://wstein.org/rh/});
  published by
  Cambridge University Press.

\item {\ptitle Algebraic Number Theory, a Computational Approach} (215 pages), (see \url{https://wstein.org/books/ant/}); under contract with the AMS.

\item {\ptitle Modular Forms, a Computational Approach}, (268
pages), published as AMS Graduate Studies in Mathematics, Volume {\bf 79},  2007,
and  available at \url{https://wstein.org/books/modform/}.

  \item {\ptitle Elementary Number Theory}
   (185 pages), published in the Springer-Verlag UTM series, 2008,
\url{https://wstein.org/ent/}.

%\item {\ptitle A Brief Introduction to Classical and Adelic Algebraic
%    Number Theory} (190 pages), \url{https://wstein.org/papers/ant/}.

\item {\ptitle Lectures on Modular Forms and Galois Representations}
  (200 pages), with K.\thinspace{}A. Ribet (on hold).


\end{books}

%\newpage
\item [Computation]
\begin{bulletlist}
\item Founder and lead developer of CoCalc
(\url{https://cocalc.com}).
\item Founder of SageMath (\url{https://www.sagemath.org}).
\item Founder and lead developer of Sage.js (\url{https://sagejs.org}),
an open mathematical system for human and agent research.
\item Fluent in TypeScript, Python, Cython, JavaScript, C/C++, LaTeX, Magma and React.
\item Wrote the modular forms, modular symbols, and modular abelian
  varieties components of Magma (over 25,000 lines of code).
%\item The Modular Forms Database: \url{https://wstein.org/tables/}.
\end{bulletlist}

\item [Selected Teaching]

%{\bf Professional Conference Organization}
%\begin{bulletlist}
%\item
%\end{bulletlist}

% {\bf Student Conference Organization}
% \begin{bulletlist}
% \item {\em Arizona Winter School}, 2007--2009, co-organizer (over 100
%   attendees each year).
% \item {\em Computing with Modular Forms}, MSRI Summer Workshop,
% July 31--August 11, 2006.  Graduate student
% workshop (about 35 attendees).
% \end{bulletlist}

{\bf University of Washington}
\setcounter{totalnumclasses}{40}
\begin{classes}
\item {\em Math 480: Sage -- Open Source Mathematical Software}, Spring 2016.
\item {\em Math 582: Computational number theory}, Winter 2016.
\item {\em Math 480: Sage -- Free Open Source Math. Software}, Spring 2014.
\item {\em Math 480: Undergraduate Number Theory}, Winter 2014.
\item {\em Math 581f: Topics in Computational Number Theory}, Fall 2013.
\item {\em Math 480a/582: Sage -- Free Open Source Math. Software}, Spring 2013.
\item {\em Math 308: Linear Algebra}, Spring 2013.
\item {\em Math 581e: Algebraic Number Theory}, Fall 2012.
\item {\em Math 480a/582: Sage -- Free Open Source Math. Software}, Winter 2012.
\item {\em Math 581g: Lectures on Modular Forms and Hecke Operators}, Fall 2011.
\item {\em Math 480a: Sage -- Free Open Source Mathematical Software}, Spring 2011.
\item {\em Math 581b: Algebraic number theory graduate course}, Fall 2010.
\item {\em Math 581d: Computer Programming for Mathematicians}, Fall 2010.
\item {\em Math 480: Computer Programming for Mathematicians}, Spring 2010.
\item {\em Math 582e: Galois Cohomomogy}, Winter 2010.
\item {\em Math 414: Elementary Number Theory}, Winter 2010.
\item {\em Math 583e: Graduate Computational Number Theory, part 2}, Spring 2009.
\item {\em Math 480: Open Source Mathematical Software},  Spring 2009.
\item {\em Math 582e: Graduate Computational Number Theory}, Winter 2009.
\item {\em SIMUW: Mathematical Finance}, Summer 2008.
\item {\em Math 480: Open Source Mathematical Software},  Spring 2008.
\item {\em Math 581f: Graduate Algebraic Number Theory}, Fall 2007.
\item {\em SIMUW: The Riemann Hypothesis}, Summer 2007.
\item {\em Math 583: The Birch and Swinnerton-Dyer Conjecture}, Spring 2007.
\item {\em Math 480: Elementary Number Theory},  Spring 2007.
\item {\em SIMUW: The Congruent Number Problem}, Summer 2006.
\item {\em Math 583: Computing with Modular Forms},  Spring 2006.
\end{classes}


\setcounter{totalnumclasses}{13}
{\bf UC San Diego}
\begin{classes}
\item {\em Elliptic Curves and Modular Forms}, Fall 2005.
\item {\em Calculus For Scientists and Engineers}, Winter 2006.
\end{classes}

\setcounter{totalnumclasses}{11}
{\bf Harvard University}
\begin{classes}
\item {\em Freshman Seminar on Fermat's Last Theorem}, Fall 2004.
\item {\em Computing With Modular Forms}, Fall 2004.
\item {\em Algebraic Number Theory}, Spring 2004.
\item {\em Modular Abelian Varieties}, Fall 2003.
\item {\em Freshman Seminar on Elliptic Curves},  Spring 2003.
\item {\em Elementary Number Theory}, Fall 2001 and Fall 2002.
\item {\em Linear Algebra}, Fall 2001 and Spring 2002.
\end{classes}
%\vspace{1ex}

%\vspace{1ex}
\setcounter{totalnumclasses}{4}
{\bf University of California at Berkeley}
\begin{classes}
 \item {\em Discrete Mathematics}, Summer 1997.
 \item {\em Calculus}, Fall 1995--Spring 1997, teaching assistant.
\end{classes}

%\vspace{1ex}
\setcounter{totalnumclasses}{2}
{\bf Northern Arizona University}
\begin{classes}
 \item {\em College Mathematics With Applications}, Spring 1995.
 \item {\em College Algebra}, Fall 1994.
\end{classes}





%{\bf Referee}
%\begin{bulletlist}
%\item Refereed papers for {\sl Mathematics of Computation},
% {\sl Mathematische Annalen}, and {\sl The Journal of Symbolic Computation}.
%\end{bulletlist}



%\item [Memberships]
%\begin{bulletlist}
%\item American Mathematical Society, 1995--present.
%\item Mathematical Association of America, 1998--present.
%\end{bulletlist}

\item [Ph.D.\\Students]

\begin{students}
\item {\bf Kevin Lui}, Ph.D., June 2019.
\item {\bf Gerardo Zelaya}, Ph.D., June 2019.
\item {\bf Hao Chen}, Ph.D. June 2016 on {\em Modular points on elliptic curves}. At Microsoft Research.
\item {\bf Simon Spicer}, Ph.D. June 2015 on {\em The.
Explicit Formula}. At Facebook.
\item {\bf Alyson Deines}, Ph.D. June 2014 on {\em Discriminant Twins}.  At CCR.
\item {\bf Robert Bradshaw}, Ph.D. June 2010 on {\em Provable
    Computation of Motivic $L$-function}.  At Google.
\item {\bf Robert Miller}, Ph.D. received June 2010 on {\em
    Verification of the Birch and Swinnerton-Dyer conjecture for
    individual elliptic curves}.  At Google.
\end{students}



\item [Other\\Activities]
{\bf Workshop and Conference Organization:}  I (co-)organized
all of the following workshops and conferences.
\begin{workshops}
%\begin{bulletlist}
\item {\em Sage Days 70}, Nov 8-14, 2015, Berkeley.
\item {\em Sage Days 68: Bug Days}, August 21-27, 2015.
\item {\em Sage Days 61: Quaternion Orders and Brandt Modules}, August 25-29, 2014, Copenhagen, Denmark.
\item {\em Sage Days 48: Notebook Dev}, June 17-21, 2013, at UW.
\item {\em Sage Days 46: Computational number theory}, February 26-March 2, 2013, Hawaii.
\item {\em Reproducibility in Computational and Experimental Mathematics}, at ICERM (Brown University), Dec 10-14, 2012.
\item {\em SIMUW: Deep Conjectures in Number Theory}, Summer 2012.
\item {\em AMS Arithmetic Statistics School}, June 24-30, 2012, in Snowbird Utah.
\item {\em Sage Days 41: The Sage Notebook}, June 11-15, 2012, at UW.
\item {\em Sage Days 40.5: Bug Days}, May 24--29, 2012 in Gold Bar, WA.
\item {\em Sage Days 36.5: Overconvergent Modular Symbols}, April 17-22, 2012, at UW.
\item {\em Sage Days 36: p-adics in Sage}, Feb 2012, San Diego, CA.
\item {\em AMS Short Course: Computing with Elliptic Curves using Sage} at the Joint Math Meetings in 2012 in Boston.
\item {\em Sage Days 32: Bug Days}, at UW.
\item {\em REU: Elliptic Curves}, 8 weeks of summer 2011, at UW.
\item {\em Sage Days 31: The Sage Notebook}, June 2011, UW.
\item {\em Sage Days 29}, March 2011, UW.
\item {\em MSRI Program in Arithmetic Statistics}, Spring 2011, at MSRI in Berkeley.
\item {\em Sage Days 27: Bug Days}, January 2011, UW.
\item {\em Sage Days 26: Women in Sage}, December 2010, UW.
\item {\em Workshop on Elliptic Curves and Computation}, October 2010, Microsoft Research.
\item {\em Sage Days 25: Numerical computation}, August, 2010, in Mubmai, India.
\item {\em Sage Days 24: Symbolic computation}, July, 2010 at RISC in Linz, Austria.
\item {\em Sage Days 23: Number theory}, July, 2010 in Leiden, Netherlands.
\item {\em Sage Days 22: MSRI Summer Graduate Student Workshop on Elliptic Curves}, June 2010 in
MSRI (Berkeley, CA).
\item {\em Sage Days 21: Function fields}, May 2010, UW.
\item {\em Sage Days 19: Bug Smash}, January 2010, UW.
\item {\em Sage Days 18: Computations related to the Birch and Swinnerton-Dyer Conjecture},
  Dec 2009, at the Clay Mathematics Institute in Cambridge, MA.
\item {\em Sage Days 17: Computing with Modular forms and L-functions}, Sep. 2009, on Lopez Island.
\item {\em Sage Days 16: Computational Number Theory}, June 2009, in Barcelona, Spain.
\item {\em Sage Days 15}, May 2009.
\item {\em Arizona Winter School: Quadratic Forms}, March 2009.
\item {\em Sage Days 14: Sage and Macaulay2 for Algebraic Geometry Experimentation}, March 2009,
MSRI (Berkeley).
\item {\em Sage Days 13: Quadratic Forms and Lattices}, March 2009, Athens, Georgia.
\item {\em Sage Days 12: Bug Smash}, Jan. 2009, San Diego, CA.
\item {\em Sage Days 11: Special functions and computational number theory meet scientific computing}, Nov. 2008, Austin, Texas.
\item {\em Sage Days 9: Mathematical graphics and visualization}, Aug. 2009, Vancouver.
\item {\em Workshop on $L$-functions and Modular Forms}, June 2008, UW.
\item {\em $L$-functions Summer School and Coding Sprint}, June 2008, UW.
\item {\em Sage Developer Coding Days}, June 2008, UW.
\item {\em Arizona Winter School: Special Functions and Transcendence },
    March 2008, Univ. of Arizona.
\item {\em Sage Days 8: Number Theory and High Performance Numerical
    Computation}, March 2008, at UT Austin.
\item {\em Sage Days 7,} Feb 2008, IPAM (UCLA).
\item {\em SAGE Days 6,} Nov. 2007, Bristol, UK.
\item {\em SAGE Days 5 -- Computational Arithmetic Geometry,} Oct 2007 at the Clay Math Institute
in Cambridge, MA.
\item {\em Workshop on Modular Forms and L-functions,} Aug. 2007 at AIM (Palo Alto).
\item {\em Sage Days 4,} June 2007 at UW.
\item {\em Modular Forms: Arithmetic and Computation,} June 2007 at Banff.
\item {\em Arizona Winter School: $p$-adic Geometry,} March 2007 at Univ. of Arizona.
\item {\em Sage Days 3,} Feb. 2007 at IPAM (UCLA).
\item {\em Interactive Parallel Computation in Support of Research in Algebra, Geometry and Number Theory,} Feb 2007 at MSRI (Berkeley).
\item {\em Sage Days 2,} Oct. 2006 at UW.
\item {\em Summer Graduate Workshop on Computing with Modular Forms,} July 2006 at MSRI (Berkeley).
\item {\em Sage Days 1,} Feb 2006 at UC San Diego.
\end{workshops}
%\end{bulletlist}


% \begin{center}
% ------------------------------------------------------------------------------------
% \end{center}

\item [Personal]
\address{}\\\newline
Phone: \phone\\
Email: \email\\
Web: \www \\
Citizenship: USA\\


%\newpage
% \item [References]\mbox{}\\
% \newcommand{\tw}{.4\textwidth}

% \hspace{-2em}\begin{tabular}{ll}
% \begin{minipage}[b]{\tw}
% {\bf Professor Benedict H. Gross}\\
% Department of Mathematics\\
% Harvard University\\
% Cambridge, MA  02138\\
% (617) 495-2172\\
% {\tt gross@math.harvard.edu}
% \end{minipage}
% &
% \begin{minipage}[b]{\tw}
% {\bf Professor Joe Harris}\\
% Department of Mathematics\\
% Harvard University\\
% Cambridge, MA  02138\\
% (617) 495-2172 \\
% {\tt harris@math.harvard.edu}
% \end{minipage}
% \\&\\
% \begin{minipage}[b]{\tw}
% {\bf Professor Hendrik W.~Lenstra}\\
% Mathematisch Instituut\\
% Universiteit Leiden\\
% The Netherlands\\
% (31/0) 71 527-7127\\
% {\tt hwl@math.leidenuniv.nl}
% \end{minipage}
% &
% \begin{minipage}[b]{\tw}
% {\bf Professor Barry C.~Mazur}\\
% Department of Mathematics\\
% Harvard University\\
% Cambridge, MA  02138\\
% (617) 495-2171 ext.~512\\
% {\tt mazur@math.harvard.edu}
% \end{minipage}
% \\&\\
% \begin{minipage}[b]{\tw}
% {\bf Professor Kenneth A.~Ribet}\\
% Department of Mathematics\\
% University of California\\
% Berkeley, CA 94720-3840\\
% (510) 642-0648\\
% {\tt ribet@math.berkeley.edu}
% \end{minipage}
% &
% \begin{minipage}[b]{\tw}
% {\bf Professor John Tate (Emeritus)}\\
% Department of Mathematics\\
% Harvard University\\
% Cambridge, MA  02138\\
% {\tt jtate@math.harvard.edu}\\
% \end{minipage}
% \end{tabular}

\end{MainList}

\end{document}





%%%%%%%%%%%%%%%%%%%%%%%%%%%%%%%%%5












\comment{
\begin{minipage}[b]{\tw}
{\bf Professor Shing-Tung Yau}\\
Department of Mathematics\\
Harvard University\\
Cambridge, MA  02138\\
(617) 495-0836\\
{\tt yau@math.harvard.edu}\\
(teaching)\\
\end{minipage}
%\begin{minipage}[b]{\tw}
%{\bf Professor Hendrik W.~Lenstra}\\
%Department of Mathematics \#3840\\
%University of California\\
%Berkeley, CA 94720-3840\\
%(510) 643-7857\\
%{\tt hwl@math.berkeley.edu}
%\mbox{}\\
%\end{minipage}

\item [References]\mbox{}\vspace{2ex}\\
\newcommand{\tw}{.4\textwidth}
\hspace{-4em}\begin{tabular}{ll}
\begin{minipage}[b]{\tw}
{\bf Professor Benedict Gross}\\
Department of Mathematics\\
Harvard University\\
One Oxford Street\\
Cambridge, MA  02138\\
(617) 495-2172\\
{\tt gross@math.harvard.edu}
\end{minipage}\\
&
\begin{minipage}[b]{\tw}
{\bf Professor Barry C.~Mazur}\\
Department of Mathematics\\
Harvard University\\
One Oxford Street\\
Cambridge, MA  02138\\
(617) 495-2171 ext.~512\\
{\tt mazur@math.harvard.edu}
\end{minipage}
\\
\begin{minipage}[b]{\tw}
{\bf Professor Kenneth A.~Ribet}\\
Department of Mathematics \#3840\\
University of California\\
Berkeley, CA 94720-3840\\
(510) 642-0648\\
{\tt ribet@math.berkeley.edu}
\end{minipage}
\end{tabular}
}

\end{MainList}

\end{document}

\begin{minipage}[b]{\tw}
{\bf Professor Hendrik W.~Lenstra}\\
Mathematisch Instituut\\
Universiteit Leiden\\
The Netherlands\\
(31/0)71 527-7127\\
{\tt hwl@math.leidenuniv.nl}\\
(research)\\
\end{minipage}


\comment{\begin{minipage}[b]{\tw}
{\bf Professor Robert F.~Coleman}\\
Department of Mathematics\\
University of California\\
Berkeley, CA 94720-3840\\
(510) 642-5101\\
{\tt coleman@math.berkeley.edu}\\
(research)\\
\end{minipage}
&\begin{minipage}[b]{\tw}
{\bf Professor Benedict Gross}\\
Department of Mathematics\\
Harvard University\\
Cambridge, MA  02138\\
(617) 495-2172\\
{\tt gross@math.harvard.edu}\\
(research)\\
\end{minipage}
\\&\\
}
\begin{minipage}[b]{\tw}
{\bf Doctor Kevin M.~Buzzard}\\
Department of Mathematics\\
Huxley Building\\
Imperial College\\
180 Queen's Gate\\
London, SW7 2BZ\\
England\\
+44 207 594 8523\\
{\tt buzzard@ic.ac.uk}
\end{minipage}
&
\begin{minipage}[b]{\tw}
{\bf Professor Robert F.~Coleman}\\
Department of Mathematics \#3840\\
University of California\\
Berkeley, CA 94720-3840\\
(510) 642-5101\\
{\tt coleman@math.berkeley.edu}\\
\mbox{}\\
\mbox{}\\
\end{minipage}
\\
&\\

%sagemathcloud={"zoom_width":110}
